\chapter{The vision pipeline, measured on the real cameras}
\label{ch:vision}

\emph{Measured on four physical cameras, 2026-08-30 -- together with
\cref{sec:hardware-session}, the only hardware numbers in this report.}

The pipeline is instrumented at thirty-three stages per camera; every stage
draws itself on the frame it was given, states its formula, and reports its
quantities with units, so a claim can be checked rather than taken on trust.
A stage has three possible outcomes, not two: it succeeds, it \emph{refuses}
and names its cause and remedy, or it is \emph{absent by design} (the
webcam's fifteen depth stages, because it has no depth sensor) -- rendering
an absence identically to a defect has cost this project findings before.

\section{The orientation defect}
\label{sec:orientation}

Three places in the software independently decided whether the scene depth
camera is mounted upside down, and disagreed. Settled by correlating each
frame's grey signature against the HD webcam (rotated by nothing, on the
same side of the room) in both orientations: the un-rotated frame agrees at
$r=\numrange{0.66}{0.70}$ over four captures, the rotated one at
$r=\numrange{0.27}{0.34}$ -- the same verdict every time, a factor of 2.4
apart. \textbf{The camera is mounted the right way up}, and every earlier
scene figure was inverted. Because the wrong rotation had been applied
\emph{consistently} to the image and the intrinsics together, the recovered
point cloud was an exact reflection of the true one and every distance,
extent and residual was preserved -- the defect was invisible to every
metric the pipeline itself computes and was found only by comparing one
camera's picture against another's.

\section{What the cameras measured}

\begin{table}[htbp]
  \centering
  \small
  \caption[Geometry recovered from each depth camera]{One frame per camera.
  The inlier fraction decides whether the row below it means anything.}
  \label{tab:vision-geometry}
  \begin{tabular}{@{}lccc@{}}
    \toprule
    & left wrist & right wrist & scene depth \\
    \midrule
    pixels returning a depth & \SI{89.4}{\percent} & \SI{80.5}{\percent} & \SI{14.5}{\percent} \\
    plane inliers & \SI{74.0}{\percent} & \SI{52.8}{\percent} & \SI{3.8}{\percent} \\
    plane residual, RMS & \SI{1.43}{\milli\metre} & \SI{1.59}{\milli\metre} & \SI{3.41}{\milli\metre} \\
    objects posed in 3D & 5 & 12 & 12 \\
    graspable within the jaw & 0 & 0 & 0 \\
    \bottomrule
  \end{tabular}
\end{table}

\textbf{The wrist cameras measure a work surface to about
\SI{1.5}{\milli\metre}} -- within a factor of three of the same code against
a renderer, the first evidence the surface-fitting stage transfers to
hardware. \textbf{The scene camera does not identify a support plane}: its
best plane holds \SI{3.8}{\percent} of the points, so every quantity that
depends on it (height, object boxes) rests on a plane one point in
twenty-six agrees with. \textbf{No object in any frame could be grasped}:
the nearest object's narrowest axis measures \SI{144}{\milli\metre} against
an \SI{85}{\milli\metre} jaw. \textbf{The open-vocabulary detector returned
nothing on any camera} for the prompt \emph{green cube}; the classical
colour-and-geometry path found the same objects it did not.

Two defects in the first runs were software, not hardware. The pick path's
own depth working band, correct for a wrist camera, applied unchanged to a
camera across the room retained \textbf{0 of 85\,061 returns} -- the band is
now declared per source. And two stages disagreed about the same frame
because depth at $480\times270$ forward-projected into colour at
$1280\times720$ lands one sample in roughly every seven pixels, so the
above-plane mask was a dot pattern that failed a connectivity check while a
per-mask stage saw it fine; splatting each projected sample raised coverage
to \SI{64.4}{\percent} and \SI{58.8}{\percent}.

The arms were to be parked at the pick pose for this session; measured
rather than assumed, the left arm was \SI{10.67}{\degree} and the right
\SI{13.53}{\degree} from the reference configuration, against a
\SI{2}{\degree} tolerance, so no figure from this session is captioned "at
the pick pose."

\section{The scene the safety layer reasons about}

\begin{figure}[htbp]
  \centering
  \includegraphics[width=\linewidth]{vision/scene_understanding.pdf}
  \caption[The scene read by three methods]{The room camera's frame,
  cropped to the wearer, the arms and the table, read by three methods.
  (b) The wearer by MediaPipe Pose: body mask and 33 landmarks. (c) The
  robot arms by a backdrop-difference rule, thinned to a centreline with a
  joint wherever it bends by more than \SI{25}{\degree}: five links on the
  left arm, four on the right. (d) The table and objects by YOLO-World
  with descriptive prompts; (e) the project's colour detector as a
  cross-check. (f) All of it on one frame. Every mask, landmark, box and
  score is the method's own output
  (\repo{figures/make\_scene\_understanding.py}).}
  \label{fig:scene-boxes-full}
\end{figure}

\Cref{fig:scene-boxes-full} reads the room camera's frame with three
methods, because no single detector covers the scene. The wearer is found
by MediaPipe Pose, the tracker the clearance floor of \cref{sec:safety}
consumes: one body, a segmentation mask covering \SI{6.0}{\percent} of the
crop, and 33 landmarks. The robot arms are white links on a white backdrop,
which defeated YOLO-World under every prompt tried and FastSAM's text
prompts alike, so they are found by an engineered rule: a 71-pixel median
of the grey image is the local backdrop level, pixels that differ from it
by more than 12 grey levels inside the corridor beside the torso, between
shoulder height and the table, are structure, FastSAM segments that overlap
a candidate by half are merged in, and the two largest components are the
arms. Each arm is thinned to its centreline (Zhang--Suen), the centreline
walked from the shoulder end to its farthest point, and a joint placed
wherever it bends by more than \SI{25}{\degree}, which is what separates one
link from the next: five links on the left arm (bends of 33, 58, 55 and
\SI{49}{\degree}) and four on the right (37, 43 and \SI{81}{\degree}). The
table and the objects are found by YOLO-World once the prompts describe
them (``white desk'' \num{0.82}, ``cardboard box'' \num{0.51}, ``green
cube'' \num{0.13}); the project's own colour detector cross-checks the
cube as a 544-pixel green blob and, once blobs on the body are excluded
(the mannequin's skin reads as red), nothing else. Every intermediate
result is kept under \repo{figures/vision/scene\_understanding/}.

What this does not give: no metric distance, because the HD camera has no
depth sensor, and a joint count that is what the silhouette shows rather
than the arm's seven joints -- at this pose the left arm presents five
visible bends and the right four. The shipped safety layer therefore
enforces the clearance floor against the tracked body of
\cref{sec:safety}, with the mannequin as the fallback it can only enlarge,
and never against a scene-camera detection of the arms.
